\chapter{Determinante}
\label{cap:determinante}

\textbf{Importante:} este capítulo é baseado em \cite{slide6doyuri}.

Neste capítulo, nós iremos introduzir formalmente o conceito de determinante. Este conceito costuma ser visto ainda no Ensino Médio, mas não de maneira tão formal como nós faremos aqui. Nós adotaremos uma abordagem axiomática, a partir da qual nós deduziremos todas as propriedades desta função. Na prática, o determinante é uma função que associa uma matriz quadrada a um número real, e que possui aplicações muito particulares em $\mathbb{R}^2$ e $\mathbb{R}^3$.

\section{Determinante}

Considere uma matriz quadrada $A_{n\times n}$ qualquer. A esta matriz, nós podemos associar um número, que é dado por algumas regras. Este número é o determinante, que nós definiremos a seguir:

\begin{Definition}[Determinante]
	Seja $A$ uma matriz quadrada $n\times n$, e sejam $a_1, \cdots, a_n$ as colunas de $A^T$ (ou seja, as linhas de $A$). Nós definimos o \textbf{determinante} de $A$ como uma função $\det:\mathbb{R}^{n\times n}\to\mathbb{R}$ que possui as seguintes três propriedades:
	
	$\mathbf{1)}$\textbf{Normalização:} $\det I = \det(e_1, \cdots, e_n) = 1$.
	
	$\mathbf{2)}$\textbf{Permutação:} $\det(a_1, a_2, \cdots, a_n) = -1\det(a_2, a_1, \cdots, a_n)$.
	
	$\mathbf{3)}$\textbf{Linearidade:} $\det(\lambda a_1 + b_1, a_2, \cdots, a_n) = \lambda\det(a_1, a_2, \cdots, a_n) + \det(b_1, a_2, \cdots, a_n)$.
\end{Definition}

A partir desta definição, nós podemos calcular, de maneira simples, o determinante de algumas matrizes $2\times 2$. Nós iremos ver três exemplos simples e, em seguida, encontramos uma fórmula geral para qualquer matriz $2\times 2$.

\begin{Example}
	Considere a matriz identidade $I = \begin{bmatrix}
		1 & 0 \\
		0 & 1
	\end{bmatrix}$. Pela propriedade $1$, nós temos que $\det I = 1$.
\end{Example}

\begin{Example}
	Considere a matriz de permutação $P = \begin{bmatrix}
		0 & 1 \\
		1 & 0
	\end{bmatrix}$. Sabemos que $P$ corresponde a uma permutação das linhas de $I$. Logo, pela propriedade $2$, nós temos que $\det P = -1$.
\end{Example}

\begin{Example}
	Considere a matriz $A = \begin{bmatrix}
		a & b \\
		a & b
	\end{bmatrix}$. Se nós fazemos uma permutação nas linhas de $A$, nós obtemos $A' = \begin{bmatrix}
		a & b \\
		a & b
	\end{bmatrix} = A$. Pela propriedade $2$, se $\det A = l$, segue que $l = -l$. Logo, $2l = 0$ e, portanto, $l = 0$. Logo, $\det A = 0$.
\end{Example}

Vamos agora enunciar e demonstrar uma fórmula para o determinante de qualquer matriz $2\times 2$.

\begin{Proposition}
	Seja $A = \begin{bmatrix}
		a & b \\
		c & d
	\end{bmatrix}$. Então, $\det A = ad - bc$.
\end{Proposition}

\begin{Proof}
	Sejam $a_1 = \begin{bmatrix}
		a \\
		b
	\end{bmatrix}$ e $a_2 = \begin{bmatrix}
		c \\
		d
	\end{bmatrix}$. Então, nós temos que:
	
	$$
	\det A = \det(a_1, a_2) = \det(ae_1 + be_2, ce_1 + de_2).
	$$
	Pela linearidade do determinante, segue que:
	
	\begin{align*}
		\det (A) &= a\det(e_1, ce_1 + de_2) + b\det(e_2, ce_1 + de_2) \\ &= a(c\det(e_1, e_1) + d\det(e_1, e_2)) + b(c\det(e_2, e_1) + d\det(e_2, e_2)).
	\end{align*}
	Como $\det(e_1, e_1) = 0$, $\det(e_2, e_2) = 0$, $\det(e_1, e_2) = 1$ e $\det(e_2, e_1) = -1$, segue que:
	
	$$
	\det A = a(c\cdot 0 + d\cdot 1) + b(c\cdot(-1) + d\cdot 0) = ad - bc.
	$$
	
	\qed
\end{Proof}

\section{Propriedades do Determinante}

A partir dos axiomas que definem o determinante, nós podemos encontrar várias outras propriedades. Nesta seção, nós enunciaremos e demonstraremos algumas delas.

\begin{Proposition}[Determinante da matriz de permutação]
	Seja $P_{n\times n}$ uma matriz de uma permutação. Então, $\det P = (-1)^k$, onde $k$ é o número de linhas que $P$ troca.
\end{Proposition}

\begin{Proof}
	Seja $\{i_1, i_2, \cdots, i_n\}$ uma permutação de $\{1, 2, \cdots, n\}$, com $k$ trocas de posição. Então:
	
	$$
	\det P = \det(e_{i_1}, \cdots, e_{i_n}).
	$$
	Pela propriedade $2$, segue que
	
	$$
	\det(e_{i_1}, \cdots, e_{i_n}) = (-1)^k\det I = (-1)^k.
	$$
	Logo, $\det P = (-1)^k$.
	
	\qed
\end{Proof}

\begin{Proposition}[Igualdade entre linhas]
	Seja $A_{n\times n}$ uma matriz com duas linhas iguais. Então, $\det A = 0$.
\end{Proposition}

\begin{Proof}
	Suponha, sem perda de generalidade, que a matriz $A$ é definida como:
	
	$$
	A = \begin{bmatrix}
		- & a_1^T & - \\
		- & a_1^T & - \\
		- & a_3^T & - \\
		& \vdots & \\
		- & a_n^T & -
	\end{bmatrix}
	$$
	Então, por um lado, nós temos que
	
	$$
	\det A = \det(a_1, a_1, a_3, \cdots, a_n).
	$$
	Por outro lado, pela propriedade $2$, nós podemos permutar as duas primeiras linhas, obtendo:
	
	$$
	\det A = -\det(a_1, a_1, a_3, \cdots, a_n).
	$$
	Assim sendo, nós temos que:
	
	$$
	\det A = -\det A\implies 2\det A = 0.
	$$
	Logo, $\det A = 0$.
	
	\qed
\end{Proof}

\begin{Proposition}[Invariância por eliminação]
	Seja $A$ uma matriz $n\times n$. Então, subtrair $\lambda$ vezes a linha $i$ da linha $j$ de $A$ não altera $\det A$. 
\end{Proposition}

\begin{Proof}
	Suponha que $\det A = \det(a_1, a_2, \cdots, a_i, \cdots, a_j, \cdots, a_n)$. Então, usando a propriedade $3$:
	
	\begin{align*}
		&\det(a_1, a_2, \cdots, a_i, \cdots, a_j - \lambda a_i, \cdots, a_n)\\ &= \det(a_1, a_2, \cdots, a_i, \cdots, a_j, \cdots, a_n) - \lambda\det(a_1, a_2, \cdots, a_i, \cdots, a_i, \cdots, a_n).
	\end{align*}
	Pela propriedade $2$, segue que $\det(a_1, a_2, \cdots, a_i, \cdots, a_i, \cdots, a_n) = 0$. Logo:
	
	$$
	\det(a_1, a_2, \cdots, a_i, \cdots, a_j - \lambda a_i, \cdots, a_n) = \det(a_1, a_2, \cdots, a_i, \cdots, a_j, \cdots, a_n) = \det A.
	$$
	
	\qed
\end{Proof}

\begin{Proposition}[Linha de zeros]
	Seja $A$ uma matriz $n\times n$ tal que $a_i = 0$, para algum $i\in\{1, \cdots, n\}$. Então, $\det A = 0$.
\end{Proposition}

\begin{Proof}
	Suponha que $\det A = \det(a_1, \cdots, 0, \cdots, a_n)$. Então, podemos escrever $\det A$ como:
	
	$$
	\det A = \det(a_1, \cdots, 0\cdot b, \cdots, a_n),
	$$
	em que $b$ é um vetor qualquer. Pela propriedade $3$, segue que:
	
	$$
	\det(a_1, \cdots, 0\cdot b, \cdots, a_n) = 0\cdot\det(a_1, \cdots, b, \cdots, a_n) = 0.
	$$
	Logo, $\det A = 0$.
	
	\qed
\end{Proof}

\begin{Proposition}[Determinante de uma matriz diagonal]
	Seja $D_{n\times n}$ uma matriz diagonal cujas entradas são $\lambda_1, \lambda_2, \cdots, \lambda_n$. Então, $\det D = \prod\limits_{i = 1}^{n} \lambda_i$.
\end{Proposition}

\begin{Proof}
	Nós temos que $\det D = \det(\lambda_1e_1, \lambda_2e_2, \cdots, \lambda_ne_n)$. Usando a propriedade $3$ repetidamente, nós temos que:
	
	\begin{align*}
		\det(\lambda_1e_1, \lambda_2e_2, \cdots, \lambda_ne_n) &= \lambda_1\det(e_1, \lambda_2e_2, \cdots, \lambda_ne_n)\\ &= \lambda_1\lambda_2\det(e_1, e_2, \lambda_3e_3, \cdots, \lambda_ne_n) \\ &\cdots \\ &= \lambda_1\lambda_2\cdots\lambda_n\det(e_1, e_2, \cdots, e_n)\\ &= \lambda_1\lambda_2\cdots\lambda_n.
	\end{align*}
	Logo, $\det D = \lambda_1\cdot\lambda_2\cdot\cdots\cdot\lambda_n$.
	
	\qed
\end{Proof}

\begin{Proposition}[Determinante de uma matriz triangular]
	Seja $T$ uma matriz triangular, e sejam $\lambda_1, \cdots, \lambda_n$ os elementos da sua diagonal principal. Então, $\det T = \prod\limits_{i = 1}^{n} \lambda_i$.
\end{Proposition}

\begin{Proof}
	Nós temos dois casos possíveis: o caso em que não há $0$ na diagonal principal, e o caso em que há algum $0$ na diagonal principal. 
	
	No primeiro caso, nós podemos fazer eliminação para chegar na matriz diagonal $D = \operatorname{diag}(\lambda_1, \cdots, \lambda_n)$ (nós não usamos a notação $\operatorname{diag}$ anteriormente, ela serve basicamente para dizer quais são as entradas de uma matriz diagonal). Assim sendo, dado que eliminação não altera o determinante, nós temos que
	
	$$
	\det T = \det D = \prod\limits_{i = 1}^{n}\lambda_i.
	$$
	
	No segundo caso, nós temos que o processo de eliminação nos dará uma linha de zeros. Assim sendo, nós obtemos:
	
	$$
	\det T = \prod\limits_{i = 1}^{n}\lambda_i = 0.
	$$
	
	Logo, $\det T = \prod\limits_{i = 1}^{n}\lambda_i$.
	
	\qed
\end{Proof}

\begin{Proposition}[Determinante em função dos pivôs]
	Seja $A$ uma matriz $n\times n$ qualquer, e sejam $d_1, \cdots, d_n$ os pivôs de $A$ (possivelmente nulos). Então, $\det A = \pm\prod\limits_{i = 1}^{n}d_i$.
\end{Proposition}

\begin{Proof}
	Considere a fatoração $A = PLU$. Então, segue-se imediatamente que $U = EP^TA$, onde $L^{-1}$ é a matriz elementar. Como $U$ é a matriz triangular que possui os pivôs de $A$ na diagonal, segue que $\det U = \prod\limits_{i = 1}^{n}d_i$. Por outro lado, nós temos que $E$, multiplicado à esquerda por $P^TA$, não altera o determinante de $P^TA$ (lembre-se que o processo de eliminação não altera o determinante, como nós já demonstramos, e $E$ é a matriz que realiza a eliminação. Logo, a matriz $E$ não altera o determinante). Assim sendo, nós temos que $\det(EP^TA) = \det(P^TA)$. Por sua vez, $\det(P^TA) = (-1)^k\det A$, onde $k$ é o número de trocas de linhas feitas, como já provamos anteriormente. Como $\det U = \det(EP^TA)$, segue que $\det A = \pm\prod\limits_{i = 1}^{n}d_i$.
	
	\qed
\end{Proof}

Faremos uma breve pausa para falar de um fato que segue diretamente desta proposição: se $A$ é singular, então $\det A = 0$. Isto é muito claro pois, se $A$ é singular, então pelo menos um dos seus pivôs é igual a zero, e o resultado sai naturalmente. Porém, há um resultado mais forte ainda que não demonstraremos aqui: $A$ é singular se, e somente se, $\det A = 0$. Ou seja, conhecendo o determinante de uma matriz quadrada, somos capazes de determinar se o seu posto é completo ou não. A demonstração deste fato será um dos exercícios deste capítulo.

Prosseguindo com as propriedades do determinante, vamos demonstrar a sua unicidade.

\begin{Theorem}[Unicidade do Determinante]
	A única função $f:\mathbb{M}^{n\times n}\to\mathbb{R}$ que satisfaz as propriedades $1$, $2$ e $3$ é o determinante.
\end{Theorem}

\begin{Proof}
	Seja $f:\mathbb{M}^{n\times n}\to\mathbb{R}$ uma outra função que satisfaz as propriedades $1$, $2$ e $3$. Então, $f$ também satisfaz todas as propriedades que nós provamos até aqui (de fato, todas elas se seguem dos três axiomas). Em particular, nós temos que $f(A) = \pm\prod\limits_{i = 1}^{n} d_i$, onde $d_i$ são os pivôs de $A$. Porém, $\pm\prod\limits_{i = 1}^{n} d_i = \det A$. Logo, $f(A) = \det A$.
	
	\qed
\end{Proof}

Tendo estabelecido a unicidade do determinante, nós podemos prosseguir para as outras propriedades.

\begin{Proposition}[Determinante do produto]
	Sejam $A$ e $B$ duas matrizes $n\times n$. Então, $\det AB = \det A\det B$.
\end{Proposition}

\begin{Proof}
	Suponha, sem perda de generalidade, que $\det B = 0$, e defina $f(A) =\dfrac{\det(AB)}{\det(B)}$. Nós queremos mostrar que $f(A) = \det(A)$. Para isso, nós vamos verificar que $f(A)$ satisfaz as propriedades $1$, $2$ e $3$.
	
	$\mathbf{1)}$ Queremos mostrar que $f(I) = 1$. Nós temos que
	
	$$
	f(I) = \dfrac{\det(IB)}{\det(B)} = \dfrac{\det(B)}{\det(B)} = 1.
	$$
	
	$\mathbf{2)}$ Queremos mostrar que $f(PA) = \pm f(A)$. Nós temos que
	
	$$
	f(PA) = \dfrac{\det((PA)B)}{\det(B)} = \dfrac{\det(P(AB))}{\det(B)} = (-1)^k\dfrac{\det(AB)}{\det(B)} = \pm f(A).
	$$
	
	$\mathbf{3)}$ Sejam $a_1, \cdots, a_n$ as linhas de $A$, e seja $A'$ a matriz definida como a matriz $A$ alterada na linha $i$ como $a_i + \lambda c_i$. Nós queremos mostrar que $f(A') = f(A) + \lambda f(a_1, \cdots, c_i, \cdots, a_n)$. Nós temos que:
	
	$$
	f(A') = \dfrac{\det(A'B)}{\det(B)}.
	$$
	
	Calculando $\det(A'B)$, nós obtemos:
	
	\begin{align*}
		\det(A'B) &= \det(a_1B, \cdots, (a_i + \lambda c_i)B, \cdots, a_nB) \\ &= \det(a_1B, \cdots, a_iB, \cdots, a_nB) + \lambda\det(a_1B, \cdots, c_iB, \cdots, a_nB) \\ &= \det(AB) + \lambda\det(\tilde{A}B),
	\end{align*}
	onde $\tilde{A}$ é a matriz $A$ alterada na linha $i$ como $c_i$. Assim sendo:
	
	$$
	f(A') = \dfrac{\det(AB) + \lambda\det(\tilde{A}B)}{\det(B)} = \dfrac{\det(AB)}{\det(B)} + \lambda\dfrac{\det(\tilde{A}B)}{\det(B)} = f(A) + \lambda f(\tilde{A}).
	$$
	Ou seja, $f(a_1, \cdots, a_i + \lambda c_i, \cdots, a_n) = f(a_1, \cdots, a_i, \cdots, a_n) + \lambda f(a_1, \cdots, c_i, \cdots, a_n)$, como queríamos provar. Logo, pelo Teorema da Unicidade do Determinante:
	
	$$
	\dfrac{\det(AB)}{\det(B)} = f(A) = \det(A)\implies \det(AB) = \det(A)\det(B).
	$$
	
	No caso em que $\det(B) = 0$, nós temos que $B$ é singular. Logo, existe $x\neq 0$ tal que $Bx = 0$. Multiplicando $A$ à esquerda de ambos os lados, nós temos que este mesmo $x$ satisfaz $ABx = 0$, ou seja, $(AB)x = 0$. Logo, $AB$ é singular, e $\det(AB) = 0$. Portanto, $\det(AB) = \det(A)\det(B)$. Repare que não é necessário tratar o caso $\det(A) = 0$, pois ele já sai direto da definição de $f(A)$.
	
	\qed
\end{Proof}

\begin{Proposition}[Determinante da inversa]
	Seja $A_{n\times n}$ uma matriz invertível. Então, $\det(A^{-1}) = \dfrac{1}{\det(A)}$.
\end{Proposition}

\begin{Proof}
	Nós temos que $\det(I) = \det(AA^{-1}) = 1$. Por sua vez, pela proposição anterior, é verdade que $\det(AA^{-1}) = \det(A)\det(A^{-1})$. Logo:
	
	$$
	\det(A)\det(A^{-1}) = 1\implies\det(A^{-1}) = \dfrac{1}{\det(A)}.
	$$
	
	\qed
\end{Proof}

\begin{Proposition}[Determinante da transposta]
	Seja $A$ uma matriz $n\times n$ qualquer. Então, $\det(A^T) = \det(A)$.
\end{Proposition}

\begin{Proof}
	Considere a fatoração $PA = LU$ (que sabemos por trezentos teoremas anteriores que existe e, sob certas condições, é única). Então, podemos escrever $A^TP^T = U^TL^T$. Por um lado, sabendo que $P^T = P^{-1}$, nós temos que:
	
	$$
	\det(A^TP^T) = \det(A^T)\det(P^T) = \det(A^T)\dfrac{1}{\det(P)} = \pm\det(A^T).
	$$
	Por outro lado, nós temos que:
	
	$$
	\det(U^TL^T) = \det(U^T)\det(L^T) = \left(\prod\limits_{i = 1}^{n} d_i\right) \cdot 1 = \prod\limits_{i = 1}^{n} d_i.
	$$
	Assim sendo, igualando os dois determinantes, segue que:
	
	$$
	\det(A^T) = \pm\prod\limits_{i = 1}^{n} d_i = \det(A).
	$$
	
	\qed
\end{Proof}

\begin{Proposition}[Determinante de uma matriz ortogonal]
	Seja $Q_{n\times n}$ uma matriz ortogonal. Então, $|\det(Q)| = 1$.
\end{Proposition}

\begin{Proof}
	Sabemos que, se $Q$ é ortogonal, então $Q^TQ = I$. Então, nós temos que:
	
	$$
	\det(Q^TQ) = \det(Q^T)\det(Q) = (\det(Q))^2 = 1.
	$$
	Logo, $\det(Q) = \pm 1$, o que quer dizer que $|\det(Q)| = 1$.
	
	\qed
\end{Proof}